\documentclass[11pt]{letter}

\usepackage[margin=0.75in]{geometry}
\usepackage{setspace}
\usepackage{hyperref}
\usepackage{xcolor}
\usepackage{enumitem}
\usepackage{parskip}

\hypersetup{
    colorlinks=true,
    urlcolor=black,
    linkcolor=black
}

\setlength{\parindent}{0pt}
\setstretch{1.05}
\date{}
\begin{document}

\begin{letter}{}

\opening{Dear Hiring Committee,}

I saw the job post on \textbf{Amazon.jobs} and thought I'd be a good fit for it.

I am a Master's student in Computer Science at the University of Cincinnati, graduating December 2026, looking for a role as a \textbf{Senior Software Development Engineer, Agentic AI} on Amazon Talent Solutions' Core Science and Engineering team. With production experience building GenAI research pipelines, ML system automation, and reliable cloud infrastructure, I am excited to support the team's work on ML, LLM, and econometric/statistical model deployment for Amazon's talent evaluation systems.

\textbf{Here are reasons that make me a great fit for the role:}

\begin{itemize}[leftmargin=*, itemsep=0.45em]
    \item \textbf{ML lifecycle management and LLM response quality:} As a Graduate Researcher, I designed a two stage memory-augmented LLM pipeline for Android bug reproduction, benchmarking against a GPT-4o + GroundingDINO + CLIP baseline and cutting latency 25.7x across 53 evaluation runs, direct hands-on work with model evaluation, debugging, and analysis techniques.

    \item \textbf{RAG and prompt engineering for GenAI applications:} I architected an evidence-grounded RAG pipeline over Gemini 2.5-flash for PolicyPilot, a medical policy intelligence portal, enforcing citation-backed answers to eliminate hallucination in high-stakes eligibility decisions, and designed a 17-tool MCP server exposing structured LLM tool-calling to external AI agents.

    \item \textbf{Scalable systems and distributed data handling:} I built an ML forecasting system with PySpark and Delta Lake processing a 2.3M-row hourly dataset across 263 zones, and at Siemens Industry Software I automated Teamcenter X infrastructure workflows with Python and Boto3, cutting deployment time 40\% across AWS environments.

    \item \textbf{Full software development lifecycle and mentorship:} As Quality Coach and SonarQube Specialist at Siemens, I increased automated code coverage from 35\% to 78\% by enforcing coding standards and CI/CD quality gates, and mentored 5 software engineers through structured onboarding, experience directly relevant to leading design and driving engineering excellence.
\end{itemize}

I am available for full-time work after my graduation in December 2026. As a Master's student at the University of Cincinnati and someone who values professionalism and growth, I am enthusiastic about contributing to Amazon Talent Solutions and helping shape how agentic AI supports hiring and development across Amazon's businesses.

Please find my CV attached below. I look forward to hearing from you. I would be excited to help build the systems that connect rigorous science with Amazon's talent decisions at global scale.

{\setlength{\parskip}{0pt}
Regards,\\
Tanmay Bhuskute,\\
bhuskutv@mail.uc.edu $\mid$ tanmay.v.bhuskute@gmail.com $\mid$ +1 7328583027,\\
3424 Brookline Ave, Flat 19, Cincinnati, OH 45220.
}

\end{letter}

\end{document}
